\documentclass[slidestop,serif,compress]{beamer}    %slidestop=title in up-left corner (with "slidescentered" centered); compress=navigation bars as small as possible

\usepackage[ugm]{mathdesign} %use Garamond font

%gets rid of bottom navigation bars and puts slide numbers right bottom
\setbeamertemplate{footline}[frame number]{}

%alternative footer: name, institute, page number
%\setbeamertemplate{footline}
%{%
%  \leavevmode%
%  \hbox{\begin{beamercolorbox}[wd=.5\paperwidth,ht=2.5ex,dp=1.125ex,leftskip=.3cm plus1fill,rightskip=.3cm]{author in head/foot}%
%    \usebeamerfont{author in head/foot}\insertshortauthor
%  \end{beamercolorbox}%
%  \begin{beamercolorbox}[wd=.5\paperwidth,ht=2.5ex,dp=1.125ex,leftskip=.3cm,rightskip=.3cm plus1fil]{title in head/foot}%
%    \usebeamerfont{title in head/foot}\insertshortinstitute \hfill \insertframenumber \,/\,\inserttotalframenumber
%  \end{beamercolorbox}}%
%  \vskip0pt%
%}

%gets rid of navigation symbols
\setbeamertemplate{navigation symbols}{}


%%complete themes (not recommended)
%%\usetheme{Madrid}             %"Berkeley":lhs navigation bar; 1st alternative: "Madrid" without a navigation bar but with a footer including name (institution), title, date page numbering (but then \institute{Univ. ...}); 2nd alternative: defaults without all the extras; 3rd alternative "Dresden": two bottom lines)
%%"Dresden is possibly easy to adapt to UvT colors
%instead of using a complete theme inner and outer themes can be specified in isolation
%inner theme:
\useinnertheme{rounded} %actually not necessary but nice for theorem boxes etc. but enumerate listings look a bit stupid! Alternative: default = just comment it out
%outer themes:
%\useoutertheme[subsection=false,footline=authorinstitutetitlepage]{miniframes} %other options for footer see beamer user guide p. 153


%%complete color themes: (ok but inner outer styles below are better)
%%\usecolortheme{albatross}    % background is colored (in blue) in albatross; the option [overlystylish] installs a background canva which is definitely too stylish for an economics department
%%\usecolortheme{seagull}   %for greyscale (black and white)
%%default is not bad!!!

%alternative to complete color schemes one can define an inner and outer color scheme
%inner color themes
%\usecolortheme{orchid} %deep color in box(theorem etc.) background fits with outer color scheme "whale"
\usecolortheme{rose} % less aggressive slight grey shading background (in boxes/theorems etc.) fits nicely with outer color scheme \usecolortheme{seahorse}
%outer color theme
\usecolortheme{dolphin}   %alternative to seahorse: "whale" but only if inner scheme is "orchid"
%
%\usefonttheme{structurebold}       %titles, headlines, foodlines and sidebars are bold

\usepackage{amsmath, eurosym, textcomp, graphicx, amsthm, amssymb,amsfonts, natbib, url, xcolor, tikz, subfig}
%\usetikzlibrary{trees}
\usepackage[english]{babel} %English language if Koma script classes are used (Bibliography instead of Inhaltsverzeichnis..)
\usepackage[latin1]{inputenc}
\usepackage[T1]{fontenc} %Encoding 
\usepackage{ae,aecompl}%avoids the problem of faint or hardly printable output


\usepackage{sgame} %Osborne's packages to draw games
%\usepackage{hyperref}
%\usepackage[all]{hypcap}


% \usepackage{comment}
% \newcommand{\com}{\comment}
% \newcommand{\ecom}{\endcomment}

%\usepackage{ngerman}
\setlength {\parindent}{0cm}
\newcommand {\bi}{\begin{itemize}}
\newcommand {\ei}{\end{itemize}}
\newcommand{\be}{\begin{enumerate}}
\newcommand{\ee}{\end{enumerate}}
\newcommand{\Ra}{\Rightarrow}
\newcommand{\ra}{\rightarrow}
\newcommand{\Lra}{\Leftrightarrow}
\newcommand{\del}{\partial}
\newcommand{\bea}{\begin{eqnarray*}}
\newcommand{\eea}{\end{eqnarray*}}
\newcommand{\less}{<}
\newcommand{\great}{>}
\newcommand{\vbar}{|}
\renewcommand{\Re}{\mathbb{R}}

\theoremstyle{plain}
\newtheorem{prop}{Proposition}
%\newtheorem{theorem}{Theorem}
\newtheorem{assumption}{Assumption}
%\newtheorem{lemma}{Lemma}
\newtheorem{proposition}{Proposition}
\newtheorem{observation}{Observation}
%\newtheorem{corollary}{Corollary}
\theoremstyle{definition}
\newtheorem{ex}{Example}
\newtheorem{exercise}{Exercise}
%\newtheorem{definition}{Definition}

%\setlength{\parskip}{1.5ex plus 0.5ex minus 0.5ex}
\author{Christoph Schottm\"uller}
\institute{University of Copenhagen}
\title{Game Theory \\Rationalizability\footnote{License: CC Attribution Share Alike 4.0}}
\date{}


\newcommand{\fr}{\frame[allowframebreaks]}
\newcommand{\ft}{\frametitle}

\newcommand{\backupbegin}{
   \newcounter{framenumberappendix}
   \setcounter{framenumberappendix}{\value{framenumber}}
}
\newcommand{\backupend}{
   \addtocounter{framenumberappendix}{-\value{framenumber}}
   \addtocounter{framenumber}{\value{framenumberappendix}}
}

\begin{document}
\maketitle
\section[Outline]{}
\frame{\ft{Outline}\tableofcontents}


\section{Rationalizability}
\label{sec:rationalizability}

\fr{\ft{Rationalizability}
\begin{itemize}
\item basic problem: which actions should we expect from a rational player?
  \begin{itemize}
  \item NE? how does a player guess the other player's actions in a one-shot game
  \end{itemize}
\item other solution concepts that require less!
\item  ``rationality'': each player must hold some (arbitrary?) belief about other players' actions and play a best response
\item $\Ra$ only actions that are best response to some belief will be played 
\item % as this holds for all players:
  $\Ra$ only best responses to action profiles that are themselves best responses have to be considered; or best responses to actions that are best responses to best responses etc.
\end{itemize}

\begin{ex}
  \begin{table}[h]
\centering
%\begin{tabular}{l|c|c|c}
 \begin{game}{2}{2}
       & L & R     \\
a1   &1,2   & 3,1 \\
a2   &0,1   & 2,3 
\end{game}
%\end{tabular}
%\caption{rationalizable strategies: example}
%\label{tab:rationalizable1}
\end{table}
\begin{itemize}
\item Could a rational P2 play R?
%   \begin{itemize}
%   \item R is best response to a2; if P2 beliefs that P1 plays a2, he should play R
%     \begin{itemize}
%     \item could a rational P2 belief that P1 plays a2?
%     \item only if a rational P1 can play a2
%     \item a2 is not a best response to any belief of P1 $\Ra$ rational P1 does not play a2
%     \end{itemize}
% \item belief that P1 plays a2 is not rational
%   \end{itemize}
% \item R is not rationalizable
\end{itemize}
\end{ex}
  
\begin{ex}
     \begin{table}[h]
\centering
%\begin{tabular}{l|c|c|c}
 \begin{game}{4}{3}
       & L & M   &  R\\ %\hline
a1   &5,14   & 4,-2 & 6,3  \\
a2   &3,0   & 5,5 &5,1 \\
a3   &6,1  &2,0  &2,16  \\
a4   &4,0  &4,2  &3,5
\end{game}
%\end{tabular}
%\caption{rationalizable strategies: example}
%\label{tab:rationalizable1}
\end{table}
\vspace*{-0.21cm}
\begin{itemize}
% \item ( one NE: $(a2,M)$ )
\item Could a rational P1 choose $a1$?
  \begin{itemize}
  \item If P1 beliefs that P2 plays $R$, a1 is b.r.

    \begin{itemize}
    \item $R$ is b.r. if P2 thinks that a3 will be played
      \item \hspace*{1cm} a3  is a best response to $L$
        \item  \hspace*{2cm}$L$ is a best response to a1
    \end{itemize}
  \end{itemize}
\item $a1$ is the action we started with!
\item $\Ra$ we cannot rule out that $a_1$ is chosen by a rational player though NE is $(a_2,M)$!
\end{itemize}
\end{ex}

% (OR definition 54.1 simplified and for 2 players):
% \begin{definition}
% An action $a_1\in A_1$ is rationalizable if there is a sequence of $((X_1^t,\mu_1^t),(X_2^t,\mu_2^t))_{t=1}^{\infty}$ that is \textbf{infinite} % (i.e. there is never a step where there does not exist a belief where there is no belief such that the required $a_j$ is not best response)
%   and of the following form:
%   \begin{enumerate}
%   \item there is some belief $\mu_1^1$ on $\Delta A_{2}$ such that $a_1$ is a best response to this belief
%   \item $X_1^1=\emptyset$ and $X_2^1=\{a_2: a_2\text{ is in the support of }\mu_1^1\}$
%   \item $\mu_2^2$: beliefs on $\Delta A_{1}$ such that $a_2$ in $X_2^1$ are best responses to $\mu_2^2$
%   \item $X_1^2=\{a_1: a_1\text{ in support of some }\mu_2^2\}$
%   \item \dots
%   \end{enumerate}
% \end{definition}
% \begin{itemize}
% \item note that the sequences above are exactly what we did in the example on the last slide
% \end{itemize}
% 
%} 

%\fr{\ft{Second definition of rationalizability}
% Second way to get rationalizability:
\begin{itemize}
\item starting point: we want a theory of rational decision making
\item requirement: self-enforcing % such a theory should not be self defeating, i.e. self-enforcing: 
(a player knowing the theory should not do better by deviating from it)
\item self enforcing implies that the theory has to prescribe \emph{sets of actions} (matching pennies)% this theory cannot only prescribe one action to every player (although this would be nice): In simple games like matching pennies this would not be self-enforcing
\item notation: theory recommends set of actions $Z_i$ to player $i$
\item what are rational beliefs? \\
  \begin{itemize}
  \item  probability zero on actions that are not in $Z=\times_{i\in N}Z_i$
\item any belief on $Z$ should be allowed: otherwise player is
    smarter than theory
  \end{itemize}
% If a player could rule out certain actions of other players, he would be smarter than the theory. 
\item self enforcing requires $Z\subseteq B(Z)$ (only best response actions to recommended actions are rational)
\end{itemize}
% These thoughts lead to the following definition:
\begin{definition}
  An action $a_i\in A_i$ is rationalizable if for each $j\in N$ there is a set $Z_j\subseteq A_j$ such that\\
(i) $a_i\in Z_i$\\
(ii) every action $a_j\in Z_j$ is a best response to a belief $\mu_j(a_j)$ of player $j$ where the support of $\mu_j$ is a subset of $Z_{-j}$.
\end{definition}
\begin{itemize}
\item if there are two families of sets $(Z_i^1)$ and $(Z_i^2)$ satisfying this definition, then $(Z_i^1\cup Z_i^2)$ satisfies the definition (check!)
\item set of rationalizable actions: $Z$
\end{itemize}

\begin{ex}[continued]
  \begin{center}
    \begin{game}{4}{3}
      & L & M & R\\ %\hline
      a1   &5,14   & 4,-2 & 6,3  \\
      a2   &3,0   & 5,5 &5,1 \\
      a3   &6,1  &2,0  &2,16  \\
      a4 &4,0 &4,2 &3,5
    \end{game}
  \end{center}

\begin{itemize}
\item why is $a1$ rationalizable? 
  \begin{itemize}
  \item what could the sets $Z_1$ and $Z_2$ be?
  %$Z_1=\{a_1,a_3\}$ and $Z_2=\{R,L\}$
  \end{itemize}
\end{itemize}

\end{ex}

\framebreak

\begin{exercise}
  \begin{itemize}
  \item Consider a Bertrand model with homogenous goods:
    \begin{itemize}
    \item two firms with 0 marginal costs set prices
    \item one  consumer buys from the firm with the lowest price
    \end{itemize}
What is the set of rationalizable prices?
  \end{itemize}
% $\Re_+$ as every price is a best response to 0
\end{exercise}


%\framebreak

% \begin{lemma}
%   The two definitions are equivalent.
% \end{lemma}
% \textbf{Proof. }\dots
% Say $a_1\in A_1$ is rationalizable according to the first definition. Then take $Z_1=\{a_1\}\cup (\cup_{t=1}^\infty X_1^t)$ and $Z_2=(\cup_{t=1}^\infty X_2^t)$. As $a_2\in X_2^t$ is a best response under $\mu_2^{t+1}$ whose support is $X_{1}^t$, the requirements of the second definition are met.\\
% Say $a_1\in A_1$ is rationalizable under the second definition. Then take $\mu_1^1=\mu_1(a_1)$ and $\mu_j^t(a_j)=\mu_j(a_j)$ for $t\geq2$. As all $X_j^t$ are subsets of $Z_j$ there is always a belief on $Z_{-j}$ rationalizing them. Hence, the defined sequence is infinite.
%
}


\section{Iterative elimination of strictly dominated actions}
\label{sec:iter-elim-strictly}



\fr{\ft{Iterative elimination of strictly dominated actions}
\begin{definition}
  An action $a_i\in A_i$ is \emph{strictly dominated }if there is a mixed strategy $\alpha_i$ of player $i$ such that $U(a_{-i},\alpha_i)\great u_i(a_{-i},a_i)$ for all $a_{-i}\in A_{-i}$.
\end{definition}
\begin{itemize}
\item a rational player does not want to use a strictly dominated action
\item $\Ra$ a rational player should not expect other players to use a strictly dominated action
\item Hence, a rational player should not use an action that is strictly dominated given that all other players play only actions that are not strictly dominated \dots
\end{itemize}

\begin{exercise}
     \begin{table}[h]
\centering
%\begin{tabular}{l|c|c|c}
 \begin{game}{3}{3}
       & L & M   &  R\\ %\hline
a1   &6,5  & 4,2  & 3,3  \\
a2   &3,0   & 3,5   &6,1 \\
a3   &4,1   &3,0    &4,16  
\end{game}
%\end{tabular}
\caption{elimination of strictly dominated strategies}
\label{tab:elim_strict_dom}
\end{table}
% In the game of table \ref{tab:elim_strict_dom}, action a3 is strictly dominated by the mixed strategy putting probability 1/2 on a1 and a2. If we do not consider a3, R is strictly dominated by a strategy putting equal weights on L and M. NOt considering a3 and M, a2 is strictly dominated by a1. Given this, M is strictly dominated by L. Hence, (a1,L) is the only strategy profile that survives iterated elimination of strictly dominated strategies.
Show that a3 is strictly dominated. Which actions survive iterated elimination of strictly dominated actions?
\end{exercise}

\begin{definition}
  The set $X\subseteq A$ survives \emph{iterated elimination of strictly dominated actions} if $X=\times_{j\in N}X_j$ and there is a collection $((X_j^t)_{j\in N})_{t=0}^T$ of sets such that for each $j\in N$
  \begin{itemize}
  \item $X_j^0=A_j$ and $X_j^T=X_j$
  \item $X_j^{t+1}\subseteq X_j^t$ for $t=0,\dots,T-1$
  \item every action in $X_j^t\backslash X_j^{t+1}$ is striclty dominated in the game $\langle N,(X_i^t),(u_i)\rangle$ i.e. the original game with actions restricted to $X^t$.
  \item No action in $X_j^T$ is strictly dominated in the game $\langle N,(X_i^T),(u_i)\rangle$.
  \end{itemize}
\end{definition}

Note: The set $X$ does not depend on the order of elimination!

\begin{lemma}[``never-best responses'' in finite games]
  An action $a_i\in A_i$ that is not a best response for any belief player $i$ might have is strictly dominated.
\end{lemma}
\textbf{Proof.} skipped

\begin{prop}\label{prop:rational_strict_el_dom}
  If $X=\times_{j\in N}X_j$ survives iterated elimination of strictly dominated action in a finite strategic game $\langle N,(A_i),(u_i)\rangle$, then $X_j$ is the set of player j's rationalizable actions.
\end{prop}
\textbf{Proof. }\emph{(for 2 players)}
\dots% \begin{itemize}
% \item preliminary: a best response of P1 to some $a_2\in X_2^t$  cannot be eliminated at stage $t$ (best response is not strictly dominated!)
% \item let $a_1\in A_1$ be rationalizable and $Z_1$ and $Z_2$ be as in the second definition of rationalizability
%   \begin{itemize}
%   \item no element of $Z_i$ will be eliminated at any stage of the elimination procedure as each element is a best response (to some belief over $Z_{-i}$)
%   \item as $a_1$ is a best response to some belief over $Z_2$ $\Ra$ cannot be eliminated $\Ra a_1\in X_1^T$
%   \end{itemize}
% \item let $a_1\in X_1^T$ (to show: $a_1$ is rationalizable)
%   \begin{itemize}
%   \item no element of $X_1^T$ is strictly dominated $\Ra$ $a_1$ is a best response to some action in $a_2\in X_2^T$ %(this is own lemma...)
%   \item is $a_1$ a b.r. to $a_2$ among all actions in $A_1$: yes, as $a_2\in X_2^t$, we cannot eliminate the best response to $a_2$ in any round
%   \item true for all actions in $X_j^T\,\Ra\,$ take $Z_j=X_j^T$ in second definition of rationalizability
\qed
%   \end{itemize}
% \end{itemize}

%%% First, assume $a_i\in A_i$ is rationalizable. Let $Z_i$ and $Z_j$ be as in the definition of rationalizability. Then  no member of $Z_j$ will be eliminated at any stage of the iterative elimination procedure as $Z_j$ is a best response to $Z_{-j}$. Hence, $Z_j\subseteq  X_j^t$ for any t. Consequently $a_i\in X_i$ (as it is a best response to some action profile in $(Z_j)$).\\
%%% Second we show that any member of $X_j$ is rationalizable. Since no member of $X_j$ is strictly dominated, every member has to be a best response \emph{among the members of $X_j$} to some belief on $X_{-j}$. We have to show that each member of $X_j$ is a best response among all members of $A_i$ to some belief on $X_{-j}$. Suppose to the contrary that there was an $a_j\in X_j$ which was not a best response (in $A_j$) to any belief on $X_{-j}$. As $a_j$ is a best response in $X_j^T$ there has to be a smallest $t$ where $a_j$ is a best response in $X_j^t$ to a belief $mu_j$ on $X_{-j}^t$. But this means that at round $t$ an action $b_j$ was eliminated that was a best response to $\mu_j$. This is the desired contradiction as $b_j$ can then not be strictly dominated
%%%\qed


\framebreak

\begin{ex}[Cournot and rationalizable strategies]\label{ex:rationalizable_cournot}
  \begin{itemize}
  \item 2- firm Cournot oligopoly with linear demand
    $p=max(1-a_1-a_2,0)$ 
  \item constant marginal costs $c$: profits are $u_1(a_1,a_2)=(1-a_1-a_2-c)a_i$
  \item Cournot equilibrium is $a_1=a_2=(1-c)/3$
  \end{itemize}
\end{ex}
\begin{ex}[Cournot and rationalizable strategies continued]
  \begin{enumerate}
  \item best response to $a_j$ is $a_i(a_j)=\max\left\{\frac{1-c-a_j}{2},0\right\}$
  \item rational player chooses $a_i$ between 0 and $1-c$
  \item $a_1 >(1-c)/2$ is never optimal if P2 chooses an action in $[0,1-c]$
  \item hence $a_2$ is at least the best response to $(1-c)/2$, i.e. $a_2\geq (1-c)/4$
  \item $a_1$ is not higher than best response to $(1-c)/4$, i.e. $a_1\leq 3(1-c)/8$
  \item $a_2$ is then at least $5(1-c)/16$
  \item \dots
  % \item \dots (every two steps the distance between the upper bound for $q_1$ and the lower bound for $q_2$ is a quarter of what it was before; however the upper bound is always above the lower bound, consequently convergence)
  \item both bounds converge to $(1-c)/3$!
  \end{enumerate}
% More formally, let the set of rationalizable actions be $Z$ with $inf\; Z=m$ and $sup\; Z=M$ (as the game is symmetric it has to be the same for all players). Note that for any belief the best response of player $i$ maximizes $E[(q_i-c)(1-q_j-q_i)]$ which is maximized by $\frac{1-E[q_j]-c}{2}$. Hence, $B_i(E[a_j])\in[B_i(M),B_i(m)]=[\frac{1-c-M}{2},\frac{1-c-m}{2}]$. For the set of rationalizable actions $Z$,we have $Z\subseteq B(Z)$ and therefore $m\geq \frac{1-c-M}{2}$ and $M\leq \frac{1-c-m}{2}$. This implies that only $m=M=(1-c)/3 $ is rationalizable.\footnote{Combining implies $1-c-2M\geq m\geq \frac{1-c-M}{2}$ and therefore $1-c\geq 3M$. In the same way, $1-c-2m\leq M\leq \frac{1-c-m}{2}$ and therefore $1-c\leq 3m$. Together with $m\leq M$, the result follows. } 
The Cournot equilibrium is the only outcome rational players can play!
\end{ex}

\begin{ex}[Cournot and rationalizable strategies continued]
  Formally, $(1-c)/3$ is the only rationalizable action:
  \begin{itemize}
  \item the game is symmetric $\Ra$ both players have the same set of rationalizable actions $Z_i$ (check!)
  \item call the highest rationalizable action $a^h$ and the lowest $a^l$
  \item if $j$ plays a mixed strategy, $i$'s best response is $a_i(a_j)=\frac{1-c-\mathbb{E}a_j}{2}$
  \item as $a^h$ is rationalizable, it must be a best response to some belief over $Z_i$ $\Ra$ $a^h\leq \frac{1-c-a^l}{2}$
  \item as $a^l$ is rationalizable, it must be a best response to some belief over $Z_i$ $\Ra$ $a^l\geq \frac{1-c-a^h}{2}$
  \item using the two inequalities, show that $a^h\leq (1-c)/3 $ and $a^l\geq (1-c)/3$
  \item conclusion: as $a^l\leq a^h$, we have $a^l=a^h=(1-c)/3$
  \end{itemize}
\end{ex}
}


\frame{\ft{Review Questions}
  \begin{itemize}
  \item Why might rational players not end up in a Nash equilibrium?
  \item What is a rationalizable action? Which conditions does it have to satisfy?
  \item What is a ``strictly dominated'' action? 
  \item Explain the method of iteratively eliminating strictly dominated actions.
  \item What is the relation between the set of rationalizable actions and the set of actions surviving iterative elimination of strictly dominated actions?
  \end{itemize}
reading: OR ch. 4 (or MSZ 4.5, 4.7 and 4.11)
}

\fr{\ft{Exercises}
  \begin{enumerate}
\item Think about the following: Why might rational players not play Nash equilibrium? Find an example game in which rational players play Nash equilibrium for sure. Find an example game where rational players might possibly choose an action profile that is not a Nash equilibrium.
% interesting aside: if in this exercise the demand is $3-p_i-2/p_j$ then both firms charge price equals marginal costs in equilibrium!
  \item Two firms compete by setting prices. The demand of firm $i$ is $D_i(p_i,p_j)=5-p_i-2/p_j$. Assume that both firms have marginal costs of 1. Write down each firm's profit function and calculate their best response functions.
% b.r.: $p_i(p_j)= 3-1/p_j$
    \begin{itemize}
    \item What is the lowest price a firm could set (with any belief over the other firm's price)? What is the highest?
%2 and 3
    \item Given that the other firm never sets prices lower/higher than the one in the last subquestion, what are the lowest and highest price a firm might set?
% 2.5 and 2.666
    \item continue to reason further\dots
% both have to charge $\frac{3+\sqrt{5}}{2}\sim 2.618$
   \item Which actions are rationalizable?
    \end{itemize}
 \item Is R a rationalizable action for P2 in the following game? (use only the definition of rationalizable action)

\begin{center}
  \begin{game}{2}{2}
    &L&R\\
T &2,1 & 1,1\\
B &1,3 & 0,1
  \end{game}
\end{center}

Which actions of P1 are rationalizable in the following game? Which of P2? (you can use all results from the lecture)

\begin{center}
  \begin{game}{3}{2}
    &L&R\\
T &2,1 & 1,1\\
M &4,2 & 0,2\\
B &1,3 & 3,1
  \end{game}
\end{center}
  \item *Show that only rationalizable actions are used with positive probability in a mixed strategy Nash equilibrium.
  \end{enumerate}
}

\end{document}